\documentclass[slidetop,11pt]{beamer}
%
% Ces deux lignes {\`a} d{\'e}commenter pour sortir 
% le texte en classe article
% \documentclass[class=article,11pt,a4paper]{beamer}
% \usepackage{beamerbasearticle}

% Packages pour les fran\c{c}ais
%
\usepackage[T1]{fontenc} 
\usepackage[latin1]{inputenc}
\usepackage[frenchb]{babel}
% pour un pdf lisible {\`a} l'{\'e}cran si on ne dispose pas 
% des fontes cmsuper ou lmodern
%\usepackage{lmodern}
\usepackage{aeguill}

% Pour afficher le pdf en plein ecran
% (comment{\'e} pour imprimer les transparents et pour les tests)
%\hypersetup{pdfpagemode=FullScreen}

% Red{\'e}finit la couleur de fond pour imprimer sur transparents
%\xdefinecolor{fondtexte}{rgb}{1,1,1}     % blanc

% commande differente pour les couleurs nomm{\'e}es - de base
%\colorlet{coultexte}{black} 

% -------------- Fioritures de style -------------
% Fait afficher l'ensemble du frame 
% en peu lisible (gris clair) d{\`e}s l'ouverture
\beamertemplatetransparentcovered

% Supprimer les icones de navigation (pour les transparents)
\setbeamertemplate{navigation symbols}{}

% Mettre les icones de navigation en mode vertical (pour projection)
% \setbeamertemplate{navigation symbols}[vertical]

\usetheme{Warsaw}

%------------ fin style beamer -------------------

% Faire appara{\^i}tre un sommaire avant chaque section
% \AtBeginSection[]{
%   \begin{frame}
%   \frametitle{Plan}
%   \medskip
%   %%% affiche en d{\'e}but de chaque section, les noms de sections et
%   %%% noms de sous-sections de la section en cours.
%   \small \tableofcontents[currentsection, hideothersubsections]
%   \end{frame} 
% }

% ----------- Contenu de la page de titre --------
\title{Recruter en informatique}
\subtitle{Les mauvaises id{\'e}es... et les bonnes !} 
\author{Gaby Wald}
\institute{\emph{to be defined}}
%% \date{\today} \date{19 mai 2020} %% \date{\today} %% \date{03 Janvier 2018}
\logo{\includegraphics[height=0.5cm]{img/logo_glider.png}}

% ----------- Nom des diff{\'e}rentes parties --------

%% Mis au fur et {\`a} mesure...

% ------------------------------------------------
% -------------   D{\'e}but document   ---------------
% ------------------------------------------------
\begin{document}
%--------- {\'e}criture de la page de titre ----------
% avec la commande frame simplifi{\'e}e
\frame[plain]{\titlepage } 
%

%------------------ Sommaire ---------------

\begin{frame}
	\frametitle{Sommaire}
	\small \tableofcontents[hideallsubsections]
\end{frame} 

\section{Recrutement en informatique, une introduction}
\begin{frame}
	\frametitle{Recrutement en informatique, une introduction}
	\tableofcontents[sections=1,currentsection,subsectionstyle=show/shaded/hide] %% sectionstyle=hide/hide,
\end{frame}

\subsection{Introduction}
\begin{frame}
	\frametitle{Introduction}
	\begin{itemize}
		\item L'objectif ici est seulement de recenser une liste de pratiques ; 
		\item Enlever les clich{\'e}s RH sur les infomaticiens, d{\'e}veloppeurs, ing{\'e}nieurs en informatique... ; 
		\item Faire sauter les soucis principaux du recrutement (nombre de candidats, turn-over, refus de contact...) ;
		\item Faire {\'e}voluer les mentalit{\'e}s (divas des deux bords, sexisme...) ; 
		\item ...  
	\end{itemize}
\end{frame}

\subsection{Mauvaises id{\'e}es, pr{\'e}jug{\'e}s...}
\begin{frame}
	\frametitle{Mauvaises id{\'e}es, pr{\'e}jug{\'e}s...}
	\begin{itemize}
		\item Recrutement uniquement sur mots-clefs et ignorer le contenu ; 
		\item Que des femmes pour recruter des hommes~\newline
		("LinkedIn est un Tinder invers{\'e}. ") ; 
		\item Recrutement sur "gadget" : comp{\'e}tences, ann{\'e}es d'exp{\'e}riences... et le sens ?
		\item Recherche de "M{\^a}les Alphas C{\'e}libataires Sans Enfants" ; 
		\item R{\'e}f{\'e}rences ? Quelles r{\'e}f{\'e}rences ? Pourquoi ?
		\item Pr{\'e}sentation g{\'e}n{\'e}rique, "recrutement sur profil", "grille de salaire"... ; et au-del{\`a} ? 
		\item "Le digital : {\`a} manipuler avec doigt{\'e} !" : num{\'e}rique !
		\item Et la culture associ{\'e}e... 
	\end{itemize}
\end{frame}

\subsection{Bonnes id{\'e}es (?)}
\begin{frame}
	\frametitle{Bonnes id{\'e}es (?)}
	\begin{itemize}
		\item Instruisez-vous, renseignez-vous, oubliez vos clich{\'e}s...
		\item Communiquer directement les enjeux (missions, r{\'e}mun{\'e}ration, conditions de travail...) ;
		\item Environnement bienveillant, tol{\'e}rant, ouvert, inclusif, divers...
		\item Distinction Vie Pro et Vie Perso...
		\item Reconnaissance et assurer l'avenir : {\^e}tre correct avec tout le monde, y compris les contractuels / prestataires / stagiaires...
		\item Une {\'e}quipe, cela se construit, et il n'y a pas d'argent magique, ni de management magique...
		%% \item ...
	\end{itemize}
	
	\emph{Note : une obligation l{\'e}gale sociale n'est pas un avantage de l'entreprise (mutuelle, transports, cong{\'e}s et RTT...). } 
\end{frame}

%%%%% %%%%% %%%%% %%%%% %%%%% 

\section{Recrutement en informatique : les mauvaises pratiques}
\begin{frame}
	\frametitle{Recrutement en informatique : les mauvaises pratiques}
	\tableofcontents[sections=2,currentsection,subsectionstyle=show/shaded/hide] %% sectionstyle=hide/hide,
\end{frame}

\subsection{Pr{\'e}jug{\'e}s "techniques"}
\begin{frame}
	\frametitle{Pr{\'e}jug{\'e}s "techniques"}
	\begin{itemize}
		\item Aspects techniques peu attirants... \c{C}a d{\'e}pend pour qui ! ; 
		\item Milieu Geek / Nerd / "Alternatif" / LGBTQIA+ / (?)... ; %% Mascu ?
		\item Terminologie magique ("digital / num{\'e}rique", acronymes, "Java / JavaScript", "Shell / PowerShell", le terminal "Matrix"...) ; 
		\item Informatique magique ("c'est simple", parfois non, apprentissage et dipl{\^o}mes !) ;  
		\item[] 
		\item Comme tout, \c{c}a s'apprend ! Au moins une culture minimale sur les acronymes, des {\'e}l{\'e}ments de pop-culture... 
	\end{itemize}
\end{frame}

\subsection{Soucis de recrutement}
\begin{frame}
	\frametitle{Soucis de recrutement}
	\begin{itemize}
		\item Turn-Over, concurrence des recruteurs (ESN / SS2I), harc{\`e}lement de contact... 
		\item Recrutement uniquement sur mots-clefs et ignorer le contenu ; 
		\item Que des femmes pour recruter des hommes, ou presque !~\newline
		("LinkedIn est un Tinder invers{\'e}. ") ; 
		\item Recrutement sur "gadget" : comp{\'e}tences, ann{\'e}es d'exp{\'e}riences... et le sens ?
		\item Recherche de "M{\^a}les Alphas C{\'e}libataires Sans Enfants" ; ...
	\end{itemize}
\end{frame}

\subsection{Lev{\'e}es de drapeaux rouges !}
\begin{frame}
	\frametitle{Lev{\'e}es de drapeaux rouges !}
	\begin{itemize}
		\item Soci{\'e}t{\'e} de Service en Informatique ("A*" / externalisations) ; 
		\item Pr{\'e}sentation comme avantages ce qui rel{\`e}ve du l{\'e}gal (cong{\'e}s pay{\'e}s, tickets restaurants, remboursement des transports...) ; 
		\item Ambiance(s) toxique(s) : management inadapt{\'e}, "mascus et alphas", prestataires / consultants...
		\item Recrutement sur profil (bloquer la concurrence ?) ; 
		\item Recrutement en comp{\'e}tition (fa\c{c}on consultant commercial) ; 
		\item Conditions : OpenSpace avec plus de 10 personnes, T{\'e}l{\'e}Travail limit{\'e}, "argent magique"...
		\item ... 
	\end{itemize}
\end{frame}

\subsection{Mauvais Cas Pratique 1}
\begin{frame}
	\frametitle{Mauvais Cas Pratique 1 : la chaine des commerciaux}
	\begin{itemize}
		%% https://www.linkedin.com/posts/rudy-onfroy_message-re%C3%A7u-sur-linkedin-votre-profil-activity-7200746078479024128-e8Vn
		\item \textbf{\emph{Exp{\'e}rience v{\'e}cue, ne reproduisez pas chez vous !}}
		\item "Votre profil correspond parfaitement {\`a} un poste que nous avons {\`a} pourvoir. J'aimerais beaucoup discuter avec vous de cette superbe opportunit{\'e}."
		\item Oui, pourquoi pas, on peut prendre RDV pour que j'en apprenne plus sur votre entreprise et sur la mission.
		\item \underline{Pendant l'entretien :} "Pourquoi avez-vous d{\'e}cid{\'e} de postuler chez nous ? Qu'est-ce qui vous int{\'e}resse dans notre entreprise ? Pourquoi pensez-vous {\^e}tes la meilleure personne pour ce poste ?"
		\item \underline{R{\'e}ponse du candidat :} "C'est vous qui m'avez contact{\'e}, je ne connaissais m{\^e}me pas votre existence hier !"
	\end{itemize}
\end{frame}

\subsection{Mauvais Cas Pratique 2}
\begin{frame}
	\frametitle{Mauvais Cas Pratique 2 : Tout dans la m{\^e}me soci{\'e}t{\'e} !}
	\begin{itemize}
		%% https://www.linkedin.com/posts/pb42_on-a-quune-chance-de-faire-bonne-impression-activity-7199329366139424768-fXHX
		\item \textbf{\emph{Exp{\'e}rience v{\'e}cue, ne reproduisez pas chez vous !}}
		\item[] 
		\item "On n'a qu'une chance de faire bonne impression,~\newline c{\^o}t{\'e} recruteur !"
		\item[] 
		\item Entretiens {\`a} r{\'e}p{\'e}tition (occuper le terrain ?) ; 
		\item Tests d{\'e}pass{\'e}s : technique d{\'e}su{\`e}te ; 
		\item Tests sans rapport avec le poste propos{\'e} ;  
		\item D{\'e}marche "{\`a} l'instinct" plut{\^o}t que de documentation ; 
		\item ... 
	\end{itemize}
\end{frame}

%%%%% %%%%% %%%%% %%%%% %%%%% 

\section{Recrutement en informatique : pour faire mieux}
\begin{frame}
	\frametitle{Recrutement en informatique : pour faire mieux}
	\tableofcontents[sections=3,currentsection,subsectionstyle=show/shaded/hide] %% sectionstyle=hide/hide,
\end{frame}

\subsection{Bonnes id{\'e}es (1)}
\begin{frame}
	\frametitle{Bonnes id{\'e}es (1)}
	\begin{itemize}
		\item Instruisez-vous, renseignez-vous : culture ? technicit{\'e} ?...
		\item Communiquer directement les enjeux (missions, r{\'e}mun{\'e}ration, conditions de travail...) ;
		\item Environnement bienveillant, tol{\'e}rant, ouvert, inclusif, diversit{\'e}...
		\item Distinction Vie Pro et Vie Perso...
		\item Reconnaissance et assurer l'avenir : {\^e}tre correct avec tout le monde, y compris les contractuels / prestataires / stagiaires...
		\item Une {\'e}quipe, cela se construit, et il n'y a pas d'argent magique, ni de management magique...
		\item Aller chercher les candidats (et s'int{\'e}resser {\`a} leur CV, leur parcours pro)...
	\end{itemize} 
\end{frame}

\subsection{Bonnes id{\'e}es (2)}
\begin{frame}
	\frametitle{Bonnes id{\'e}es (2)}
	\begin{itemize}
		\item Poste sur lequel il y a recrutement : les enjeux doivent {\^e}tre donn{\'e}s (premier contact ou entretien qui suit) ;~\newline
		Pr{\'e}parez ces donn{\'e}es en amont : sinon, c'est du temps perdu pour tout le monde !
		
		\item Environnement de travail inclusif, et diversit{\'e}... ~\newline
		Pas de concessions ici : mal {\`a} l'aise en entretien ou apr{\`e}s, le ressenti a des cons{\'e}quences pas forc{\'e}ments visibles imm{\'e}diatement.
		
		\item Au-del{\`a} de la passion : des individus !~\newline
		{\'E}vitez quand m{\^e}me les divas et / ou les personnes qui se n{\'e}gligent (franchise et direct l{\`a}-dessus si le cas se pr{\'e}sente). 
	\end{itemize} 
\end{frame}

\subsection{Bonnes id{\'e}es (3)}
\begin{frame}
	\frametitle{Bonnes id{\'e}es (3)}
	\begin{itemize}
		\item Faites propositions concr{\`e}tes sur profils susceptibles de r{\'e}pondre avec les crit{\`e}res suivants de proposition :
		\begin{itemize}
			\item Mission int{\'e}ressante (pr{\'e}ciser le contexte autant que possible, t{\^a}ches {\`a} accomplir...) ; 
			\item Adaptation au profil (exemple : {\'e}viter de proposer des t{\^a}ches routini{\`e}res {\`a} quiconque a un bac+5 ou bac+2 avec plus de 3 ans d'exp{\'e}riences) ; 
			\item Indiquez TOUJOURS un salaire, ou une fourchette de salaire (renseignez-vous !), {\`a} partir du budget ; 
			\item {\'E}ventuellement "on s'aligne !" ; 
			\item Cibler des personnes qui ont {\'e}t{\'e} embauch{\'e}es en sortie d'{\'e}cole 
				et qui sont en poste depuis au moins 2-3 ans (sous-pay{\'e}es) ;
			\item Cibler personne d{\'e}j{\`a} connue favorablement : ad{\'e}quation, projets pr{\'e}c{\'e}dents, voire culture interne connue ; 
		\end{itemize} 
		\item ...
	\end{itemize} 
\end{frame}

%%%%% %%%%% %%%%% %%%%% %%%%% 

\section{Autres id{\'e}es et r{\'e}f{\'e}rences}
\begin{frame}
	\frametitle{Autres id{\'e}es et r{\'e}f{\'e}rences}
	\tableofcontents[sections=4,currentsection,subsectionstyle=show/shaded/hide] %% sectionstyle=hide/hide,
\end{frame}

\subsection{Autres id{\'e}es (En vrac)}
\begin{frame}
	\frametitle{Autres id{\'e}es (En vrac - 1/8)}
	%% --- 
	\begin{columns}[T]
	\begin{column}[T]{5cm}
		\includegraphics[width=5cm]{img/1651166248397.jpeg}~\\
	\end{column}
	\begin{column}[T]{7cm}
		\begin{itemize}
			\item Tester le processus de recrutement !
			\item Candidat myst{\`e}re : point de vue ext{\'e}rieur. 
			\item Pourquoi le candidat refuse ou ne r{\'e}pond plus ?
			\item Priorit{\'e}s du candidat ? (poste, localisation, r{\'e}mun{\'e}ration...)
			\item L'ordre des priorit{\'e}s est important !
			\item Ressenti d'int{\'e}r{\^e}t mutuel : le CV est-il lu avec attention ? 
		\end{itemize}
	\end{column}
	\end{columns}
	%% \begin{itemize}
	%%	\item Tester le processus de recrutement !~\newline
	%%		\texttt{https://www.linkedin.com/posts/activity-6925493203215245312-cR3U}~\newline
	%%		L'art de la guerre...~\newline
	%%		Parfois, au vus de certaines tentatives en mati{\`e}re de \#recrutement : cr{\'e}er une soci{\'e}t{\'e} de conseil pour le recrutement dans la tech ; un \#concept simple, des candidats myst{\`e}res se glissent dans votre processus de recrutement... \#startup \#conseil \#tech \#processus \#candidatmyst{\`e}re~\newline
	%%		De la m{\^e}me fa\c{c}on que des clients myst{\`e}res vont dans des magasins, restaurants, services... Et font ensuite un compte-rendu de fa\c{c}on anonyme, en g{\'e}n{\'e}ral via leur soci{\'e}t{\'e} ou entreprise (genre Gault et Millau, Guide du Routard, Guide Michelin...). Oui, cela existe encore ! \#GAFAM~\newline
	%%		Oui, de nombreuses entreprises auraient besoin de ce retour d'exp{\'e}rience de la part des candidats re\c{c}us ou contact{\'e}s pendant un recrutement, surtout dans "la tech" ! Et cela compte pour leur r{\'e}putation...~\newline
	%%		\#num{\'e}rique \#digital \#entreprises \#entreprise \#r{\'e}flexion
	%% \end{itemize} 
\end{frame}

\begin{frame}
	\frametitle{Autres id{\'e}es (En vrac - 2/8)}
	\begin{itemize}
		\item "Le recrutement, c'est comme la chasse : il y a les bons, les brutes et les truands !" (Le recrutement fonctionne mieux avec les bons !)
		\item Il ne faut pas que ce soit \emph{The Most Dangerous Game}, \emph{Battle Royale} ou \emph{Hunger Games} ! (Il y a des codes sociaux accept{\'e}s et des limites l{\'e}gales. )
		\item Demandez ce que les candidats pensent des recruteurs ! ("fraichement sorti d'{\'e}cole de commerce" ; "stagiaire" ;~\newline "recruteuse usant de (ses) charmes" ; "chasseur de pigeons" ;~\newline "la professionnelle", ... )
	\end{itemize}
	%% \begin{itemize}
	%% 	\item \texttt{https://www.linkedin.com/posts/activity-6932730981455949825-Tqwx}~\newline
	%% 		En mati{\`e}re de \#Recrutement, je songe {\`a} {\'e}crire un \#livre ou un \#JdR (au moins parodique) : il y a tellement {\`a} exprimer sur le sujet !~\newline
	%% 		Le commercial de "SSII / ESN du CAC 40" (oui, celle-l{\`a}) qui a probablement fait une {\'e}cole de commerce ; le stagiaire qui d{\'e}barque en RH...~\newline
	%% 		(le commercial d'{\'e}cole de commerce, le stagiaire qui d{\'e}barque, la "recruteuse" convaincue de ses charmes, le vendeur qui recherche des pigeons en demandant des r{\'e}f{\'e}rences, la RH professionnelle, le recruteur trop honn{\^e}te...).
	%% \end{itemize} 
\end{frame}

\begin{frame}
	\frametitle{Autres id{\'e}es (En vrac - 3/8)}
	%% ---
	\begin{columns}[T]
	\begin{column}[T]{5cm}
		\includegraphics[width=5cm]{img/1655474560502.jpeg}~\\
	\end{column}
	\begin{column}[T]{7cm}
		~\\ \hfill ~\\
		La pissoti{\`e}re : le d{\'e}rangement doit valoir le co{\^u}t ! Sinon, c'est g{\^e}nant...~\\~\\~\\
		
		Laissez vos coordonn{\'e}es ; la relance c'est quand le processus de recrutement est d{\'e}j{\`a} entam{\'e} !
		%% ~\\ \hfill ~\\
	\end{column}
	\end{columns}
	%% \begin{itemize}
	%%	\item \texttt{https://www.linkedin.com/posts/activity-6943563576917807104-djRZ}~\newline
	%%		La pissoti{\`e}re : le d{\'e}rangement doit valoir le co{\^u}t ! Sinon, c'est g{\^e}nant...
	%% \end{itemize} 
\end{frame}

\begin{frame}
	\frametitle{Autres id{\'e}es (En vrac - 4/8)}
	\begin{columns}[T]
	\begin{column}[T]{5cm}
		\includegraphics[width=5cm]{img/1657616562690.jpeg}~\\
	\end{column}
	\begin{column}[T]{7cm}
		\begin{itemize}
			\item Jusqu'o{\`u} aller pour recruter ? 
			\item La p{\^e}che au chalut ? La drague ? La prostitution ?
			\item {\`A} part pour le porno (films et produits d{\'e}riv{\'e}s) : utilit{\'e} du genre dans le recrutement ?
		\end{itemize}%% ~\\
		\includegraphics[width=5cm]{img/1657616562426.jpeg}~\\
	\end{column}
	\end{columns}
	%% \begin{itemize}
	%%	\item \texttt{https://www.linkedin.com/posts/activity-6952547782712803328-t7Om}~\newline
	%%		On le rappelle : LinkedIn n'est pas Tinder ! Ah ben si en fait !
	%% \end{itemize} 
\end{frame}

\begin{frame}
	\frametitle{Autres id{\'e}es (En vrac - 5/8)}
	\begin{itemize}
		\item \textbf{Personnalisez les messages !}
		\item En soi, faire des messages g{\'e}n{\'e}rique est, peut-{\^e}tre, une bonne id{\'e}e ; 
			sauf pour les destinataires, qui n'auront aucune information (mission, poste, r{\'e}mun{\'e}ration...). 
		\item \textbf{Mod{\`e}le pr{\'e}-rempli ou {\`a} remplir : }
		\begin{itemize}
			\item Pourquoi cibler telle ou telle personne ?
			\item Descriptif de poste lisible rapidement !
			\item R{\'e}mun{\'e}ration ?!
		\end{itemize}
		%% \item \texttt{https://www.linkedin.com/posts/activity-6963075115613876224-uY\_J}~\newline
		%% 	Arr{\^e}tez les messages g{\'e}n{\'e}riques !~\newline
		%% 	Les messages de \#recrutement g{\'e}n{\'e}riques sur \#LinkedIn : tout un art ! Entre autre ici avec Extia~\newline
		%% 	C'est comme les lettres de \#motivation : le mod{\`e}le \#g{\'e}n{\'e}rique ou \#template vaut le co{\^u}t d'aller plus loin !~\newline
		%% 	Pourquoi mon profil ? Descriptif de poste ? R{\'e}mun{\'e}ration ? ...~\newline
		%% 	Idem Atos : les messages de recrutement g{\'e}n{\'e}riques ne sont pas appr{\'e}ci{\'e}s. Pour de bonnes raisons.~\newline
		%% 	\#Recrutement \#Num{\'e}rique '\#Digital' \#LinkedIn \#BonnesPratiques ...~\newline
		%% 	Qui r{\'e}pond encore ? Je n'aime ni le "name dropping" ni le "name shaming" : que faire contre ?~\newline
		%% 	Poke Devoteam \#recruteurs : pouvez vous mieux former vos recruteuses et recruteurs. Les messages g{\'e}n{\'e}riques de \#Recrutement \#IT : ce n'est pas en votre faveur ! Surtout en double ce matin.~\newline
		%% 	\#ESN \#SS2I \#LinkedIn ...~\newline
		%% 	Pour les autres soci{\'e}t{\'e}s de services aussi !
	\end{itemize} 
\end{frame}

\begin{frame}
	\frametitle{Autres id{\'e}es (En vrac - 6/8)}
	\begin{columns}[T]
	\begin{column}[T]{5cm}
		\includegraphics[width=5cm]{img/1673900833331.jpeg}~\\
	\end{column}
	\begin{column}[T]{7cm}
		\begin{itemize}
			\item Le contenu du poste ?
			\item Et sa localisation ?
			\item Et la r{\'e}mun{\'e}ration ? (ni la visibilit{\'e} ni la passion)
		\end{itemize}%% ~\\
	\end{column}
	\end{columns}
	%% \begin{itemize}
	%%	\item \texttt{https://www.linkedin.com/posts/activity-7020848968561520641-lC2k}~\newline
	%%		\#Twitter : monde merveilleux o{\`u}, {\`a} l'image de la vraie vie, on peut trouver des critiques de \#LinkedIn et des m{\'e}thodes de \#Recrutement foireuses !~\newline
	%%		\#IRL \#R{\'e}seauxSociaux ...~\newline
	%%		Pr{\'e}sentez le contenu du poste ! Et poste r{\'e}aliste !
	%% \end{itemize} 
\end{frame}

\begin{frame}
	\frametitle{Autres id{\'e}es (En vrac - 7/8)}
	\begin{columns}[T]
	\begin{column}[T]{5cm}
		\includegraphics[width=5cm]{img/1687203978939.jpeg}~\\
	\end{column}
	\begin{column}[T]{8cm}
		\begin{itemize}
			\item Le recrutement "efficace" de non-candidats pour les nuls.~\newline
			\item Lui proposer un poste sans d{\'e}crire la stack. 
			\item Lui proposer un poste {\`a} 200km. 
			\item Lui proposer un poste sous pay{\'e}.~\newline
			\item Perdu, rien de tout \c{c}a ne marche ! 
		\end{itemize}%% ~\\
	\end{column}
	\end{columns}
	%% \begin{itemize}
	%%	\item https://www.linkedin.com/posts/activity-7076646406790246400-c2Ac~\newline
	%%		Le recrutement "efficace" de non-candidats pour les nuls.~\newline
	%%		Lui proposer un poste sans d{\'e}crire la stack.~\newline
	%%		Lui proposer un poste {\`a} 200km.~\newline
	%%		Lui proposer un poste sous pay{\'e}.~\newline
	%%		Perdu, rien de tout \c{c}a ne marche dans la r{\'e}alit{\'e}. \#recrutement \#dev \#emploi \#travail \#informatique .~\newline
	%% \end{itemize} 
\end{frame}

\begin{frame}
	\frametitle{Autres id{\'e}es (En vrac - 8/8)}
	\begin{itemize}
		\item Int{\'e}ressez-vous {\`a} la culture des personn{\'e}es cibl{\'e}es ! 
		\item \c{C}a ne vous dit rien ? Heureusement, il y a \emph{Wikipedia}, \emph{WikiHow}, \emph{Wookiepedia}...
		\item Sinon, les moteurs de recherches peuvent r{\'e}pondre {\`a} vos questions, les personnes cibl{\'e}es aussi ; 
		\item Gothique, Chasseur ou Chasseuse de Pok{\'e}mon, Gameuse, R{\^o}liste, Intello ? ... 
		\item Et sinon, vous buvez quoi ? (r{\'e}ponses possibles : eau, th{\'e}, caf{\'e}, soda, autre). 
	\end{itemize}
	%% \begin{itemize}
	%% 	\item \texttt{https://www.linkedin.com/posts/activity-6980604583660118016--eer~}\newline
	%% 		Recrutement en \#informatique : \c{c}a se passe mieux quand il y a connaissance de la culture et des fa\c{c}ons de faire du \#num{\'e}rique !~\newline
	%% 		Am{\'e}lioration continue ? Culture ?~\newline
	%% 		\#RH \#PopCulture '\#digital' ...
	%% \end{itemize} 
\end{frame}

%% \def\sectionPartBibliographie{Bibliographie / Mediagraphie}
%% \section{\sectionPartBibliographie}
%% \begin{frame}
%% 	\frametitle{\sectionPartBibliographie}
%% 	{\footnotesize 
%% 		\nocite{*}
%% 		%toutes references biblio : 6 lettres + 2 chiffres
%% 		\bibliography{presentationRecrutementInformatique}
%% 		% \bibliographystyle{frplain} % plain or frplain
%% 		\bibliographystyle{plain}
%% 	}
%% \end{frame}

\end{document}
